\clearpage
\section{ПРАКТИЧЕСКАЯ РЕАЛИЗАЦИЯ ПРОЕКТА GAMIFIED TASK CRM}

\subsection{Общая структура реализации и конфигурация среды}

Практическая реализация проекта выполнена в репозитории \textit{gamified\_app} и ориентирована на web-режим запуска Flutter-приложения. Входной точкой системы является файл \texttt{lib/main.dart}, в котором организованы базовые этапы инициализации: подготовка Flutter binding, инициализация локального хранилища Hive, регистрация адаптеров моделей, запуск сервиса аутентификации и подъем корневого виджета приложения в \texttt{ProviderScope}. Такое построение старта соответствует требованию предсказуемого bootstrap-процесса и исключает скрытые неявные зависимости между модулями \cite{repoReadme}.

С точки зрения конфигурации окружения проект использует параметры \texttt{--dart-define}. К обязательным для backend-first контура относятся \texttt{SUPABASE\_URL} и \texttt{SUPABASE\_ANON\_KEY}; для AI-функций подключается \texttt{MINIMAX\_API\_KEY}. Дополнительно предусмотрены флаги режимов: \texttt{DEMO\_MODE=true} и \texttt{AI\_DEMO\_MODE=true}. Эта схема важна как с эксплуатационной, так и с безопасностной точки зрения: секреты не хранятся в исходном коде, а режимы исполнения разделяются явно \cite{repoReadme,owaspM4}.

Маршрутизация организована через \texttt{GoRouter}. В приложении заданы два главных маршрута: \texttt{/login} и \texttt{/app}. Логика \texttt{redirect} обеспечивает автоматическое перенаправление неаутентифицированного пользователя на экран входа и блокирует доступ в закрытую часть интерфейса без валидной сессии. При этом для авторизованного пользователя обратный переход на \texttt{/login} также блокируется. Такой механизм устраняет класс ошибок, когда часть защищенного функционала случайно остается доступной по прямому URL \cite{goRouterDocs,owaspM1}.

Пользовательская оболочка реализована в \texttt{RootLayout} как \texttt{IndexedStack} с пятью вкладками: главная панель, задачи, CRM, отчеты и профиль. Индексная модель хранения экранов позволяет сохранять локальное состояние вкладок и ускоряет переключение между ними, что особенно заметно при частой навигации между задачами и клиентскими карточками. Дополнительные действия создания сущностей (задача, клиент) вынесены в контекстные формы через \texttt{FloatingActionButton} и \texttt{ModalBottomSheet}, что делает поток работы компактным и однотипным во всех основных модулях.

В практической части курсового проекта важно подчеркнуть, что архитектурные решения выбирались исходя из учебной цели: показать контролируемую, расширяемую систему, а не точечный прототип. Поэтому структура проекта поддерживает функциональную декомпозицию по доменным блокам:
\begin{itemize}
  \item \texttt{lib/features/dashboard/presentation/};
  \item \texttt{lib/features/tasks/presentation/};
  \item \texttt{lib/features/crm/presentation/};
  \item \texttt{lib/features/reports/presentation/};
  \item \texttt{lib/features/profile/presentation/};
  \item \texttt{lib/features/ai/presentation/}.
\end{itemize}

Такое деление улучшает поддерживаемость: разработчик работает в рамках конкретной предметной зоны, не затрагивая нерелевантные части системы. Кроме того, модульная структура облегчает контроль требований при защите курсового проекта, поскольку каждую функциональную область можно объяснить отдельно: назначение, входные данные, основной сценарий, ограничения и тестовый контур.

\subsection{Модель данных и backend-first интеграция}

Практическая реализация данных строится на двух взаимосвязанных источниках: локальное хранилище (demo-режим) и облачный backend (рабочий режим). На уровне клиентского кода модель задачи описывается классом \texttt{Task}, включающим поля \texttt{id}, \texttt{title}, \texttt{description}, \texttt{priority}, \texttt{category}, \texttt{clientId}, \texttt{dueDate}, \texttt{createdAt}, \texttt{isCompleted}. Для CRM используется модель \texttt{Client} с идентификатором, контактными атрибутами, статусом, источником и датой создания.

Для backend-контура в README проекта зафиксирован минимальный состав таблиц: \texttt{public.tasks}, \texttt{public.crm\_clients}, \texttt{public.user\_stats}. Ключевой миграцией выступает добавление связи \texttt{tasks.client\_id} с таблицей клиентов, что обеспечивает трассируемость сценария \textit{задача, привязанная к конкретному клиенту} \cite{repoReadme,migrationClientId,bootstrapSchema}.

С инженерной точки зрения важна не только структура таблиц, но и дисциплина применения схемы. В проекте предусмотрен идемпотентный SQL-скрипт bootstrap, создающий/обновляющий базовые таблицы, внешние ключи, индексы и триггеры \texttt{updated\_at}. Такой подход снижает риск дрейфа схемы между окружениями и упрощает повторную настройку стенда перед защитой.

В таблице \ref{tab:data-model} приведена укрупненная характеристика ключевых сущностей системы.

\begin{table}[h!]
\caption{Ключевые сущности и их роль в системе}
\centering
\begin{tabular}{|p{3.5cm}|p{5.2cm}|p{5.5cm}|}
\hline
\textbf{Сущность} & \textbf{Основные поля} & \textbf{Назначение в бизнес-потоке} \\
\hline
Task & id, title, priority, status, due\_date, client\_id & Планирование и контроль операционных действий пользователя, связь с клиентом при необходимости \\
\hline
CRM Client & id, first\_name, last\_name, email, phone, status, source & Хранение клиентской базы и контекста взаимодействия, сегментация контактов \\
\hline
User Stats & xp, level, streak, last\_task\_date & Геймифицированные метрики активности и мотивации пользователя \\
\hline
\end{tabular}
\label{tab:data-model}
\end{table}

Как видно из таблицы \ref{tab:data-model}, модель данных не перегружена второстепенными сущностями. Это осознанное решение, направленное на устойчивость базового сценария. Слишком раннее усложнение схемы (например, введение десятков справочников) в учебном проекте обычно приводит к росту технического долга и снижению качества тестирования.

Важной частью реализации является обработка ошибок backend-взаимодействия. В \texttt{main.dart} реализован метод нормализации исключений в пользовательские сообщения. Примеры покрываемых случаев: отсутствие таблиц в Supabase, ошибки политики RLS, некорректный UUID, истекший JWT-токен. Такая прослойка уменьшает когнитивную нагрузку на пользователя и ускоряет диагностику проблем во время демонстрации проекта.

Если формализовать надежность клиент-серверного обмена на уровне сессии, можно использовать показатель доступности операций:
\begin{equation}
A = \frac{N_{ok}}{N_{ok} + N_{err}},
\label{eq:availability}
\end{equation}

где $N_{ok}$ --- количество успешно завершенных операций, $N_{err}$ --- количество операций, завершившихся ошибкой. В практической разработке показатель \eqref{eq:availability} используется как оперативная оценка стабильности интеграции после изменений в схеме или политике доступа.

\subsection{Реализация пользовательских сценариев и модульной логики}

Ключевой прикладной сценарий проекта включает последовательность действий: вход пользователя в систему, создание карточки клиента, постановка задачи с привязкой к клиенту, проверка отображения результата в списке задач и в деталях клиента. Этот сценарий отражает основной смысл CRM-подхода: управление работой строится не изолированно, а в контексте конкретных контактов и сделок \cite{repoReadme,docsDefenseScript}.

В модуле задач реализованы базовые CRUD-операции, приоритеты, категории и сроки выполнения. Привязка к клиенту сделана nullable-полем, что поддерживает две бизнес-модели: задачи без клиентского контекста (внутренние операции) и задачи с клиентским контекстом (внешние взаимодействия). В CRM-модуле реализовано управление карточками клиентов и их статусами, что позволяет отслеживать жизненный цикл взаимодействия от лида до активного клиента.

Модуль отчетности в текущем контуре предназначен для представления агрегированных показателей, а не для сложной BI-аналитики. Такой уровень детализации соответствует масштабу курсового проекта: пользователь получает базовую оценку результативности, при этом система остается легкой для сопровождения. При дальнейшем развитии проекта модуль отчетов можно расширить когортным анализом, сегментацией по источникам лидов и динамикой конверсии по периодам.

Важное практическое свойство реализации --- единый UX-паттерн ввода данных через нижние модальные формы. Он снижает вероятность сценарных расхождений между модулями и ускоряет обучение пользователя работе с системой. В учебном контексте это также упрощает проведение защиты: демонстрация строится на повторяемых действиях, а не на случайных переходах по сложному интерфейсу.

В таблице \ref{tab:scenarios} приведены основные сценарии и ожидаемые результаты.

\begin{table}[h!]
\caption{Ключевые пользовательские сценарии и ожидаемый эффект}
\centering
\begin{tabular}{|p{4.5cm}|p{5.2cm}|p{4.5cm}|}
\hline
\textbf{Сценарий} & \textbf{Проверяемое поведение} & \textbf{Практический результат} \\
\hline
Авторизация пользователя & Redirect между \texttt{/login} и \texttt{/app} в зависимости от сессии & Защита закрытого контура и корректный старт работы \\
\hline
Создание клиента & Валидация полей и сохранение карточки в локальный/backend контур & Актуальная CRM-база для планирования задач \\
\hline
Создание задачи с client\_id & Сохранение связи \textit{задача-клиент} и отображение в двух представлениях & Трассируемость операционной активности \\
\hline
Просмотр отчета & Агрегация базовых показателей по активности & Быстрая обратная связь по прогрессу \\
\hline
Обработка ошибок backend & Нормализация технической ошибки в понятное сообщение & Снижение времени диагностики и доработки \\
\hline
\end{tabular}
\label{tab:scenarios}
\end{table}

Как видно из таблицы \ref{tab:scenarios}, практическая реализация охватывает полный минимальный контур пользовательской ценности: вход в систему, данные клиентов, постановка задач и контроль результата.

Отдельно следует отметить геймифицированную составляющую проекта. В модели \texttt{UserStats} учитываются показатели XP, уровень и streak. Несмотря на упрощенный характер этой подсистемы, она решает важную задачу мотивации пользователя к регулярному выполнению задач. С точки зрения продуктовой логики это оправдано: даже простые механики прогресса повышают вовлеченность в системах персональной продуктивности \cite{productivityBook}.

\subsection{Верификация качества и результаты практической проверки}

Проверка качества в практической главе организована как последовательность обязательных шагов. На первом уровне выполняется \texttt{flutter analyze} для выявления статических проблем. На втором уровне запускаются автоматизированные тесты: \texttt{flutter test} и \texttt{flutter test --coverage}. На третьем уровне применяются скрипты контроля покрытия и интеграционного smoke-потока. На четвертом уровне предусмотрен web E2E-сценарий с реальной аутентификацией и backend-операциями \cite{repoReadme}.

С методической точки зрения такая организация проверки дает несколько преимуществ:
\begin{itemize}
  \item раннее обнаружение синтаксических и типовых ошибок;
  \item контроль корректности доменной логики на уровне модулей;
  \item подтверждение целостности основного пользовательского сценария;
  \item снижение риска регрессий перед защитой проекта.
\end{itemize}

В таблице \ref{tab:quality-checks} представлена сводка контрольного контура качества.

\begin{table}[h!]
\caption{Контур проверок качества проекта}
\centering
\begin{tabular}{|p{4.2cm}|p{5.5cm}|p{4.5cm}|}
\hline
\textbf{Инструмент} & \textbf{Что проверяется} & \textbf{Назначение} \\
\hline
\texttt{flutter analyze} & Предупреждения, потенциальные ошибки типов и API-вызовов & Статическая устойчивость кода \\
\hline
\texttt{flutter test} & Unit и widget сценарии & Корректность логики и UI-компонентов \\
\hline
\texttt{flutter test --coverage} + gate скрипты & Полнота тестового охвата runtime-критичного кода & Контроль инженерной дисциплины \\
\hline
Integration smoke script & Базовый пользовательский поток в интеграционной среде & Быстрая проверка работоспособности после изменений \\
\hline
Web Supabase E2E script & Сквозной сценарий с реальным backend и авторизацией & Подтверждение эксплуатационной готовности \\
\hline
\end{tabular}
\label{tab:quality-checks}
\end{table}

Результаты практической проверки показывают, что выбранный набор инструментов является достаточным для учебного проекта, где важны не абсолютные метрики enterprise-уровня, а воспроизводимость и аргументированность качества. Наличие четко описанных команд запуска также облегчает приемку преподавателем: сценарии можно повторить без ручного вмешательства в код.

Для формализации проверки к отчету добавлены приложения: полный набор команд запуска и проверки (Приложение 1), структура ключевых директорий (Приложение 2), а также таблица приемочных критериев (Приложение 3). Ссылки на эти материалы позволяют быстро сопоставить теоретическое описание и практический контур реализации.

При этом необходимо зафиксировать ограничения. Во-первых, в проекте пока не реализован полный контур нагрузочного тестирования, поэтому поведение при высоком числе одновременных пользователей оценивается косвенно. Во-вторых, модуль отчетности сфокусирован на базовых сводках и не покрывает сложные BI-сценарии. В-третьих, AI-подсистема рассматривается как вспомогательный модуль и не является критичным контуром бизнес-логики.

Указанные ограничения не снижают учебную ценность проекта, поскольку главная цель курсовой работы состоит в демонстрации корректно спроектированного и реализованного базового решения. Более того, наличие явно сформулированных ограничений является признаком зрелого инженерного подхода, когда команда осознает текущие границы системы и планирует дальнейшее развитие.

\subsection{Экономия трудозатрат, риски и направления развития}

Практический эффект от реализации \textit{Gamified Task CRM} оценивается через сокращение ручных операций и упорядочивание взаимодействия с клиентами. В простейшем сценарии пользователь получает единое рабочее пространство вместо разрозненных таблиц, мессенджеров и заметок. Это снижает потери контекста и улучшает управляемость ежедневных задач.

С точки зрения проектных рисков можно выделить три основные группы:
\begin{enumerate}
  \item \textbf{Конфигурационные риски}: ошибки в переменных окружения, несовпадение схемы БД, некорректные политики доступа.
  \item \textbf{Продуктовые риски}: избыточная сложность интерфейса при масштабировании функционала.
  \item \textbf{Технические риски}: накопление технического долга при быстром добавлении новых модулей без рефакторинга.
\end{enumerate}

Для управления этими рисками в проекте применяются базовые меры: bootstrap-скрипт схемы, автоматизированные проверки, модульная структура и явные runtime-флаги режимов. В дальнейшем целесообразно усилить практику за счет централизованной телеметрии, мониторинга ошибок и формализации release-процесса.

Приоритетные направления развития после завершения курсового проекта могут быть сформулированы следующим образом:
\begin{itemize}
  \item расширение аналитического блока (воронка, сегментация, динамика конверсии);
  \item добавление ролевой модели пользователей и разграничения прав;
  \item внедрение уведомлений по срокам и событиям;
  \item развитие мобильного канала на базе той же кодовой базы Flutter;
  \item усиление quality-контура за счет дополнительных интеграционных и контрактных тестов.
\end{itemize}

Суммарно практическая реализация подтверждает достижение промежуточного целевого состояния: создан работоспособный, структурно устойчивый и расширяемый программный продукт, соответствующий учебным требованиям по предметной области, архитектуре, данным и качеству выполнения.
